\begin{Exercise}[label={ex:radice}]
  Scrivere un programma che dato un numero positivo ne restituisca la radice intera
\end{Exercise}

\begin{Answer}[ref={ex:radice}]
	Abbiamo già visto in un'esercitazione precedente lo pseudocodice per questo algoritmo.
	
	Iniziamo con la soluzione utilizzando il costrutto \texttt{while}
  \lstinputlisting[style=cstyle]{src/radice_intera1.c}

	Alternativamente si può utilizzare anche il ciclo \texttt{for}
	\lstinputlisting[style=cstyle]{src/radice_intera2.c}

	\paragraph{Attenzione: }
			In questa soluzione il ciclo non è vuoto, ma semplicemente l'unica operazione è il decremento della variabile che viene eseguito nella terza espressione del ciclo \texttt{for}
\end{Answer}
